\section{Speculative decoding: MTP versus DFlash2}\label{sec:spec}

Two drafters were available for this model. \textbf{MTP} is a prediction head trained with the model
and stored in the same GGUF file; it costs no extra weights but keeps its own one-layer \flag{f16}
draft cache. \textbf{DFlash2} is a separate 1.1\,GB block-diffusion drafter that proposes a block of
up to 7 tokens in one parallel pass, conditioned on hidden states tapped from five target layers
\citep{chen2026dflash,incoai2026dflash2_card}. Without speculation this machine decodes a 27B model at
12--23\tps{} into a shallow cache, depending on build and engine, and at 3.5\tps{} at 246K tokens, so a drafter is not optional for
interactive use.

\subsection{Speculation is a different decode path, not a free speed knob}\label{sec:nonequiv}
For nine days my notes said: speculative decoding is lossless under greedy decoding, so it changes
speed only, and no accuracy run needs a speculative arm. Perplexity runs for speculative
configurations were skipped on that basis. The first experiment that tested the premise refuted it.

\begin{table}[t]
\centering\small
\caption{Speculation at context 32{,}768, first battery. \Qsix, greedy (temperature 0, top-$p$ 1, fixed
seed), 1{,}024-token limit, all 164 HumanEval+ problems, zero empty responses. ``Byte-exact'' compares
each completion with the no-speculation baseline. The DFlash2 row ran on a different engine build
(E-B) and an older drafter file, so its equivalence and speed are not comparable with the MTP
rows.}\label{tab:nonequiv}
\shrinktofit{\begin{tabular}{@{}lrrrrrr@{}}
\toprule
configuration & tok/s (median) & speed-up & acceptance & byte-exact & HE\,\% & HE+\,\% \\
\midrule
no speculation & 18.46 & 1.00$\times$ & --- & 164/164\textsuperscript{a} & 94.5 & 91.5 \\
MTP $n{=}2$ & 37.44 & 2.03$\times$ & 0.954 & \textbf{131/164}\textsuperscript{a} & 93.9 & 90.2 \\
MTP $n{=}4$ & 47.03 & 2.55$\times$ & 0.892 & \textbf{131/164} & 93.9 & 90.9 \\
\emph{DFlash2 $n{=}4$ (E-B)} & \emph{51.78} & \emph{2.80$\times$} & \emph{0.917} & \emph{132/164} & \emph{93.3} & \emph{90.2} \\
\bottomrule
\multicolumn{7}{@{}l}{\footnotesize\textsuperscript{a}\,Re-run unchanged a day later: all 164 completions byte-identical to the first run, in both arms.}
\end{tabular}}
\end{table}

\Cref{tab:nonequiv}: with MTP on, 33 of 164 completions differ from the plain-decoding output, at
temperature 0 with a fixed seed. A repeat a day later reproduced \emph{both} arms byte for byte. So
this is not run-to-run noise. The engine is deterministic, and speculation deterministically takes a
different path for about one problem in five. Depth 2 and depth 4 each change 33 problems but share
only 25 of them.

\textbf{What this does and does not mean.} Task accuracy did not move: pass@1 differs by at most 1.3
points across the four rows, far inside $\pm4.6$. The changed completions are different programs, not
worse ones. ``One in five'' is also a fact about HumanEval+'s output lengths: the chance of divergence
rises with completion length, from 2.4\,\% of the shortest quartile of problems to about 50\,\% of the
longest, a hazard near one divergence per 950 generated tokens that grows with position. That pattern
fits accumulating numerical differences better than a broken acceptance rule.

\textbf{The mechanism is not established.} The leading explanation is batch-shape numerics: verifying
$n{+}1$ tokens in one pass runs different GPU kernels from decoding one token, the logits differ in
their last bits, and a near-tie flips. Kernels of this kind are deterministic run to run yet vary with
batch shape \citep{thinkingmachines2025nondeterminism}, which is exactly what byte-identical repeats
plus cross-configuration differences look like. At one point I read the determinism result as ruling
this explanation out; it does not, and that reading is withdrawn. Others report the same behaviour on
llama.cpp with quantized targets \citep{llamacpp_issue25618}, and vLLM states its guarantee as
lossless ``up to hardware numerics'' \citep{vllm_specdecode_docs}. I did not compare logits at a
divergence point, so this remains the leading hypothesis. All of it is greedy-only; sampling at
temperature above zero was not tested.

\begin{takeaway}
Treat the speculative setting as part of the configuration under test. A reproducible pipeline must
pin it, not only the seed, and a quantization comparison should run without speculation on both arms
or with the identical setting on both.
\end{takeaway}

\subsection{First-battery results, and why they cannot rank the drafters}
The DFlash2 row of \cref{tab:nonequiv} was the fastest at 32K. At depth, the first battery recorded
DFlash2 failing to load above 163{,}840 tokens and collapsing to 8\tps{} there. I reported that as
``fastest when shallow, dominated when deep'' and then had to take most of it back:
\begin{itemize}
\item the at-depth speed was timed over a \emph{17-token} generation (the probe sent a corpus to the
chat endpoint with no instruction; the model acknowledged it and stopped). The same defect voided
MTP's first at-depth figures and every historical acceptance of exactly 1.000;
\item the load failures at 213K and 262K were the drafter's 738\,MiB weight buffer failing to fit on
GPU0 under a split tuned for MTP: a placement problem, not a DFlash2 limit;
\item the arm ran on a different engine build, with a drafter file that predated an upstream
metadata fix, at draft depth 4 where the drafter's block size of 8 makes 7 the design point.
\end{itemize}
Between the batteries DFlash2 support was merged into upstream llama.cpp
\citep{llamacpp_pr27342_dflash2}, which removed the engine confound and made a fair comparison
possible.

\subsection{Same build, both drafters, to 246K tokens}\label{sec:s15speed}

\begin{figure}[t]
\centering
\begin{tikzpicture}
\begin{groupplot}[group style={group size=2 by 1,horizontal sep=1.2cm,ylabels at=edge left},
  paper,width=0.5\linewidth,height=6.2cm,
  xlabel={filled context (thousand tokens)},ylabel={decode (tok/s), median},
  xmin=0,xmax=270,ymin=0,ymax=58,xtick={32,131,246},
  legend style={font=\scriptsize}]
\nextgroupplot[title={\small\Qsix},legend to name=speclegend,legend columns=3,
  error bars/y dir=both,error bars/y explicit,error bars/error bar style={line width=0.5pt}]
\addplot[cbred,mark=triangle*] coordinates {(32.768,48.39) += (0,5.8) -= (0,13.7) (131.072,31.50) += (0,4.7) -= (0,4.7) (246.415,25.68) += (0,5.7) -= (0,5.4)};
\addplot[cbblue,mark=*] coordinates {(32.768,29.39) += (0,0.3) -= (0,4.2) (131.072,21.56) += (0,3.1) -= (0,1.7) (246.415,17.29) += (0,3.1) -= (0,2.1)};
\addplot[cbblue,mark=o,dashed] coordinates {(32.768,26.94) (131.072,19.28) (246.415,14.15)};
\addplot[cbred,mark=triangle,dashed] coordinates {(32.768,28.60) (131.072,22.89) (246.415,17.46)};
\addplot[black!65,mark=square*,densely dotted] coordinates {(32.768,11.66) (131.072,5.61) (246.415,3.49)};
\legend{DFlash2 $n{=}7$,MTP $n{=}4$,MTP $n{=}2$,DFlash2 $n{=}3$,no speculation}
\nextgroupplot[title={\small\Qfour},
  error bars/y dir=both,error bars/y explicit,error bars/error bar style={line width=0.5pt}]
\addplot[cbred,mark=triangle*] coordinates {(32.768,44.01) += (0,1.3) -= (0,5.9) (131.072,33.95) += (0,9.9) -= (0,5.0) (246.415,27.66) += (0,8.7) -= (0,4.7)};
\addplot[cbblue,mark=*] coordinates {(32.768,29.58) += (0,1.5) -= (0,3.1) (131.072,22.06) += (0,4.8) -= (0,2.7) (246.415,16.33) += (0,4.3) -= (0,3.8)};
\addplot[cbblue,mark=o,dashed] coordinates {(32.768,27.92) (131.072,20.17) (246.415,14.66)};
\addplot[cbred,mark=triangle,dashed] coordinates {(32.768,29.11) (131.072,21.54) (246.415,16.88)};
\addplot[black!65,mark=square*,densely dotted] coordinates {(32.768,13.14) (131.072,5.93) (246.415,3.61)};
\end{groupplot}
\node at ($(group c1r1.south)!0.5!(group c2r1.south)+(0,-1.75cm)$) {\pgfplotslegendfromname{speclegend}};
\end{tikzpicture}
\caption{Decode rate against filled context, both drafters on one upstream build (E-C). Greedy,
identical code prompts, \flag{q4\_0} KV, 1{,}024-token continuations, 8 per cell (4 at 32K); markers are
medians and bars the full range across continuations. At its design depth ($n{=}7$) DFlash2 leads at
every depth on both builds. At a shared shallow depth ($n{=}3$, dashed) the two drafters are
indistinguishable. Without speculation decode falls to 3.5\tps{} at 246K.}\label{fig:specdepth}
\end{figure}

The second battery ran both drafters on build E-C with the current drafter file, each at several draft
depths, each as a separate server launch. (On this build the draft depth cannot be changed per
request: the per-request fields are compiled out, while the properties endpoint still echoes a value,
so a per-request sweep would silently measure one depth.) Every cell was filled with the same real
code to 32K, 131K or 246K tokens, then asked for 1{,}024-token continuations.

\Cref{fig:specdepth} gives the result. At 246{,}415 filled tokens DFlash2 at $n{=}7$ decodes at 25.7
against MTP's best 17.3\tps{} on \Qsix{} (+49\,\%) and 27.7 against 17.0 on \Qfour{} (+62\,\%). At
131K the margins are +45\,\% and +54\,\%; at 32K, +64\,\% and +45\,\%. MTP's best depth is 4 in two of those cells and 3 in the other four (\cref{app:speed}). The ranges of
DFlash2 at $n{=}7$ and the best MTP depth overlap in one cell only (\Qsix{} at 246K, by 0.15\tps).
Three details matter for anyone reproducing this:

\textbf{The advantage appears at DFlash2's design depth and not at a shared shallow depth.} At a shared depth of 3 the two drafters are
within noise of each other everywhere. MTP gains little from going deeper (its depth 2, 3 and 4
medians are within 10--22\,\% of each other, and depth 4 does not beat depth 3 consistently), while DFlash2 gains 40--70\,\% from depth 3 to 7 because
its block is drafted in one parallel pass.

\textbf{DFlash2's accepted length grows with depth of context.} Mean tokens accepted per verification
pass rose from 4.4 to 4.6 on \Qsix{} and from 3.5 to 4.6 on \Qfour{} between 32K and 246K. The
``acceptance collapse at depth'' in my early notes was the short-generation artifact described above.
Raw acceptance ratios are not comparable between drafters (they divide by different depths); report
accepted length and tokens per second.

\textbf{The decode speed-up from speculation grew with filled context.} Against no speculation the best drafter
gives $3.4$--$4.2\times$ at 32K, $5.6$--$5.7\times$ at 131K and $7.4$--$7.7\times$ at 246K, because
each verification pass amortises an ever more expensive attention read over several tokens. DFlash2's
prefill is also 34--52\,\% faster at depth (201 against 135\tps{} at 246K on \Qsix), since MTP
prefills its draft cache as well.

The pre-registered selection rule compared decode on thinking-mode requests of about 100K input tokens
(three per configuration; median [min, max]). It chose DFlash2 on both builds: 20.4 [19.8, 22.7]
against 16.3 [16.1, 16.6]\tps{} on \Qsix, cleanly; 29.8 [23.8, 51.3] against 17.4 [15.9, 22.1] on
\Qfour, where the ranges do not overlap but one range spans a factor of two over three requests. The matched-depth probe above is the stronger evidence for the same conclusion. The
accuracy batteries in \cref{sec:aalcr} then served as a long-generation speed check over 756 real
answers: DFlash2 decoded at a median 34.6--48.9\tps{} against MTP's 23.7--28.2, at lower energy per
request (28--41\,kJ against 45--78\,kJ).

\subsection{Draft depth and the acceptance trap}
The first battery swept MTP depth $\{2,4,8\}$ at matched depths of 131K and 196K on all four builds,
three 512-token generations per cell, to learn the best depth per build. It could not answer.
Repetition-to-repetition decode spread reached 166\,\%, and acceptance within one cell ranged by up to
0.63. Two lessons survive.

\emph{Clustered events give falsely tight intervals.} Pooled over 1{,}000--4{,}000 draft events, a
Wilson interval on acceptance is $\pm0.02$. The events inside one generation are strongly correlated;
the effective sample is the three generations, and their spread is $\pm0.3$.

\emph{``Acceptance falls with draft depth'' is nearly a tautology.} It fell in 10 of 11 adjacent pairs
(one-sided sign test $p=0.006$, which fails a Holm correction at its rank). For any per-token match
probability $q<1$, the ratio accepted/drafted must fall as depth grows. Inverting the
stop-at-first-mismatch formula gives $q\approx0.68$--$0.95$, roughly constant across depth within a
cell. $q$ carries the information; the ratio does not. At the full 262K window, MTP $n{=}8$ failed to
load at all: its draft context no longer fits.

\subsection{Where decode variance comes from}\label{sec:variance}
Early speed tables on this host showed 17--41\,\% spread between identical repetitions, which I first
called unexplained noise. Two reviewers noted that decode correlates with acceptance at $r^2$ of
0.83--0.99. That correlation is an identity, not a finding: llama.cpp's speculative decode rate is
exactly $(1+n\cdot\text{acceptance})/t_\text{pass}$, since each target pass emits one guaranteed token
plus the accepted drafts. Reconstructing generated length from pass and acceptance counts recovers 190
to 191 of 192 tokens in all 47 cells. What is informative is the residual. Dividing throughput by
tokens-per-pass, with nothing fitted, collapses within-configuration spread from 16.7--40.7\,\% to
1.3--13.4\,\%; 98.7\,\% of within-cell variance is tokens-per-pass and 0.9\,\% is pass time. The
``noise'' was sampling variance of the generated text: those cells ran at temperature 0.7, so each
repetition wrote a different continuation with a different acceptance. Pass time itself is well
behaved, linear in filled depth at about 0.45\,\textmu s per token. \textbf{For controlled comparisons, run speed probes greedy with fixed prompts}, and report sampled
workload throughput separately, since real acceptance depends on the text being generated. The second
battery did both (this probe, and the 756 sampled answers of \cref{sec:aalcr}). Its medians and ranges
describe the continuations that were run within one launch per cell; they are not a significance test.
